\documentclass[11pt,a4paper]{article}

% Compile with pdfLaTeX. Replace the blue prompts and add the three
% Cinderella screenshots before submitting the final PDF.
\usepackage[margin=25mm]{geometry}
\usepackage[T1]{fontenc}
\usepackage{lmodern}
\usepackage{amsmath,amssymb}
\usepackage{booktabs,array,tabularx}
\usepackage{graphicx}
\usepackage{float}
\usepackage{xcolor}
\usepackage{hyperref}
\hypersetup{colorlinks=true,linkcolor=black,urlcolor=blue}
\setlength{\parindent}{0pt}
\setlength{\parskip}{0.55em}
\renewcommand{\arraystretch}{1.2}

\newcommand{\studentname}{\textit{[Student name]}}
\newcommand{\studentid}{\textit{[Student ID]}}
\newcommand{\prompt}[1]{\textcolor{blue!60!black}{\textbf{[To complete: #1]}}}
\newcommand{\blank}{\underline{\hspace{1.5cm}}}
\newcommand{\shortblank}{\underline{\hspace{0.8cm}}}
\newcommand{\screenshot}[1]{%
  \IfFileExists{#1}{\includegraphics[width=0.88\linewidth,height=70mm,keepaspectratio]{#1}}{%
    \fbox{\parbox[c][53mm][c]{0.83\linewidth}{\centering
      \prompt{Insert Cinderella screenshot: \texttt{\detokenize{#1}}}\\[1ex]
      Show the assembled mechanism and its DOF annotation.}}}}

\begin{document}
\hypersetup{pageanchor=false}
\begin{titlepage}
  \centering
  \vspace*{32mm}
  {\Large Kinematics of Mechanisms\par}
  \vspace{9mm}
  {\huge\bfseries Homework 1\par}
  \vspace{14mm}
  {\Large Report and Solutions\par}
  \vfill
  \begin{tabular}{@{}ll@{}}
    Student name: & \studentname\\[2mm]
    Student ID: & \studentid\\[2mm]
    Submission date: & \textit{[Date]}
  \end{tabular}
  \vspace*{20mm}
\end{titlepage}

\hypersetup{pageanchor=true}
\tableofcontents
\newpage

\section{Theoretical Questions}

\subsection{Question 1: Mechanism and machine}
A \emph{mechanism} is an arrangement of rigid bodies (links) connected by
kinematic pairs so that their relative motion is constrained. Its purpose is
to transmit or transform motion, such as converting a crank's continuous
rotation into the oscillation of a rocker. A four-bar linkage is a mechanism.

A \emph{machine} is an assembly that uses mechanisms together with other
components to perform a useful task by transmitting or converting energy and
forces. A reciprocating engine, for example, contains a slider-crank
mechanism that converts piston motion into shaft rotation, while the complete
engine also includes the energy source and supporting systems. Thus,
``mechanism'' emphasizes the prescribed relative motion; ``machine''
emphasizes useful function and power transfer. A machine may contain several
mechanisms, and a mechanism can be part of a machine.

\subsection{Question 2: Four-bar linkage and its types}
A planar four-bar linkage is a closed chain of four rigid links joined by four
revolute (pin) joints. One link is fixed as the frame; the other three are
commonly called the input link, coupler, and output link. For an ordinary
nondegenerate linkage, the planar mobility equation gives
\begin{equation}
  M=3(n-1)-2J_1-J_2=3(4-1)-2(4)-0=1,
  \label{eq:fourbar-mobility}
\end{equation}
where $n$ includes the frame, $J_1$ counts one-DOF joints, and $J_2$
counts two-DOF joints. One independent input therefore fixes the
configuration within a chosen assembly branch.

The usual types describe how the two links attached to the frame move:
\begin{itemize}
  \item \textbf{Double-crank (drag link):} both adjacent links can make full
    revolutions relative to the frame.
  \item \textbf{Crank-rocker:} one adjacent link can make a full revolution;
    the other only oscillates through a bounded angle.
  \item \textbf{Double-rocker:} neither adjacent link can make a full
    revolution; both oscillate.
\end{itemize}
Which type results depends on both the link lengths and the choice of fixed
link. The coupler's motion is not itself used to name these three types.
The strict Grashof classifications are given in Question~4.

\subsection{Question 3: Main lower kinematic pairs and their DOF}
A lower kinematic pair connects two bodies through contacting surfaces.
Its DOF is the number of independent \emph{relative} motion coordinates
allowed by the joint. In three-dimensional space, the common lower pairs
are as follows.

\begin{table}[H]
  \centering
  \caption{Main lower pairs and relative spatial mobility.}
  \begin{tabularx}{\linewidth}{@{}l c X@{}}
    \toprule
    Pair & DOF & Permitted relative motion\\
    \midrule
    Revolute (R) & 1 & Rotation about a fixed axis.\\
    Prismatic (P) & 1 & Translation along a fixed axis.\\
    Helical (H) & 1 & Coupled rotation and translation along an axis, with a fixed screw pitch.\\
    Cylindrical (C) & 2 & Independent rotation about and translation along the same axis.\\
    Spherical (S) & 3 & Rotation about three independent axes through a common center; no translation.\\
    Planar (E) & 3 & Two independent translations in a plane and one rotation normal to it.\\
    \bottomrule
  \end{tabularx}
\end{table}

The values in Table~1 are \emph{spatial} joint DOFs. For the planar mobility
count in Part~A, a planar revolute or prismatic pair is a one-DOF joint,
whereas a general planar higher pair permits two relative DOFs. Joint DOF
must not be confused with the mobility of the complete mechanism.

\subsection{Question 4: Grashof, non-Grashof, and special case}
Let $S$ and $L$ be the shortest and longest lengths of a four-bar chain,
and let $P$ and $Q$ be the other two lengths. For a nondegenerate four-bar,
Grashof's criterion compares the two sums:
\begin{equation}
  \begin{array}{rcll}
    S+L &<& P+Q & \text{strict Grashof chain},\\
    S+L &>& P+Q & \text{non-Grashof chain},\\
    S+L &=& P+Q & \text{Grashof equality, or change-point case}.
  \end{array}
  \label{eq:grashof}
\end{equation}

For a \textbf{strict Grashof chain}, at least the shortest link can rotate
fully relative to another link. Fixing the shortest link gives a
double-crank; fixing a link adjacent to the shortest gives a crank-rocker;
fixing the link opposite the shortest gives a double-rocker. These statements
assume distinct ordinary link lengths and a feasible assembly.

For a \textbf{non-Grashof chain}, no link can rotate through a complete
revolution relative to the other links. With any link fixed, the ordinary
inversion is a double-rocker.

At \textbf{equality}, $S+L=P+Q$, the linkage can reach a collinear
\emph{change-point} configuration. Its motion there is singular: the usual
classification becomes a limiting case, and continuation through the
aligned position can depend on the assembly branch or physical guidance.
It should therefore be reported separately instead of being treated as an
ordinary strict Grashof or non-Grashof linkage.

\section{Part A: Cinderella Models and Mobility}
Create each mechanism in Cinderella, label its mobility with the
\texttt{Add text} tool, and insert a screenshot below. Submit the three
corresponding \texttt{.cdy} files with the report. For each model, identify
\emph{rigid bodies}, including the ground, before counting joints. If more
than two links meet at one pin, count the equivalent binary joints
appropriately. Explain any pin-in-slot or contact modeling convention.

For a general planar mechanism, use the Gr\"ubler--Kutzbach count
\begin{equation}
  M=3(n-1)-2J_1-J_2,
  \label{eq:mobility}
\end{equation}
where $n$ is the number of links including ground, $J_1$ the number of
one-DOF lower pairs, and $J_2$ the number of two-DOF higher pairs.
Check the count against observed independent motions, especially if a
mechanism has redundant constraints or a special geometry.

\subsection{Case 1: Compound linkage}
The dimensions supplied in the assignment are
\begin{equation*}
\begin{gathered}
  AB=2,\ BC=9,\ CD=6,\ DE=6,\ EQ=16,\ QG=6,\\
  EF=8,\ QL=8,\ DP=9,\ PG=7,\ PA=4,\\
  \angle DEF=15^\circ,\qquad \angle LQG=15^\circ.
\end{gathered}
\end{equation*}
\begin{figure}[htbp]
  \centering\screenshot{figures/part_a_case1.png}
  \caption{Case 1 Cinderella model with its DOF annotation.}
\end{figure}

\prompt{List the rigid links and identify all pin joints, including any
compound joints.} The count is
\begin{align*}
  n&=\blank,\qquad J_1=\blank,\qquad J_2=\blank,\\
  M&=3(n-1)-2J_1-J_2=\blank.
\end{align*}
\prompt{Explain any special geometric or redundant constraints and verify
the observed mobility in Cinderella.}

\subsection{Case 2: Casement window mechanism}
Use $AB=4$, $BC=8$, $CD=5$, $DE=7$, $EF=5$, with the fixed angles
$\theta=125^\circ$ and $\varphi=125^\circ$. Points $D$ and $F$ are sliders.
\begin{figure}[htbp]
  \centering\screenshot{figures/part_a_case2.png}
  \caption{Case 2 Cinderella model with both sliders and its DOF annotation.}
\end{figure}

\prompt{Identify each rigid link, each revolute pair, and each prismatic
pair; explain how the fixed angles are represented.}
\begin{align*}
  n&=\blank,\qquad J_1=\blank,\qquad J_2=\blank,\\
  M&=3(n-1)-2J_1-J_2=\blank.
\end{align*}
\prompt{Verify the observed mobility in Cinderella.}

\subsection{Case 3: Bull gear and slotted bar}
Use $AC=4$, $AB=8$, $AP=8$, and $KB=20$.
\begin{figure}[htbp]
  \centering\screenshot{figures/part_a_case3.png}
  \caption{Case 3 Cinderella model, including gear, slot, slider, and DOF annotation.}
\end{figure}

\prompt{Describe the constraints at the gear, connecting rod, slot, and
horizontal slider. State whether a pin-in-slot is counted as a higher pair
or modeled with an equivalent slider and lower pairs; use one consistent
count.}
\begin{align*}
  n&=\blank,\qquad J_1=\blank,\qquad J_2=\blank,\\
  M&=3(n-1)-2J_1-J_2=\blank.
\end{align*}
\prompt{Verify the observed mobility in Cinderella.}

\begin{table}[htbp]
  \centering
  \caption{Part A mobility results.}
  \begin{tabular}{@{}lcccc@{}}
    \toprule
    Mechanism & Links $n$ & One-DOF joints $J_1$ & Two-DOF joints $J_2$ & Mobility $M$\\
    \midrule
    Case 1 & \blank & \blank & \blank & \blank\\
    Case 2 & \blank & \blank & \blank & \blank\\
    Case 3 & \blank & \blank & \blank & \blank\\
    \bottomrule
  \end{tabular}
\end{table}

\section{Part B: Grashof Condition and Linkage Classification}
\subsection{Task 1: Classify three four-bar chains}
All lengths in Table~\ref{tab:partb-input} are in centimetres. In each case,
sort the four lengths into $S,L,P,Q$, compare $S+L$ with $P+Q$, and name
the mechanism types obtainable by selecting different fixed links.

\begin{table}[htbp]
  \centering
  \caption{Given four-bar link lengths in centimetres.}
  \label{tab:partb-input}
  \begin{tabular}{@{}lrrrr@{}}
    \toprule
    Case & $l_1$ & $l_2$ & $l_3$ & $l_4$\\
    \midrule
    1 & 12 & 4 & 10 & 9\\
    2 & 4 & 8 & 11 & 7\\
    3 & 6 & 3 & 5 & 4\\
    \bottomrule
  \end{tabular}
\end{table}

\begin{table}[htbp]
  \centering
  \caption{Part B, Task 1 working table.}
  \begin{tabularx}{\linewidth}{@{}lcccccX@{}}
    \toprule
    Case & $S$ & $L$ & $P,Q$ & $S+L$ & $P+Q$ & Class and possible fixed-link types\\
    \midrule
    1 & \shortblank & \shortblank & \shortblank & \shortblank & \shortblank & \prompt{Class and types}\\
    2 & \shortblank & \shortblank & \shortblank & \shortblank & \shortblank & \prompt{Class and types}\\
    3 & \shortblank & \shortblank & \shortblank & \shortblank & \shortblank & \prompt{Class and types}\\
    \bottomrule
  \end{tabularx}
\end{table}

\prompt{For each case, show the inequality or equality explicitly. State
which link is fixed for every claimed double-crank, crank-rocker, or
double-rocker inversion. Treat a change-point case separately.}

\subsection{Task 2: Design a double-crank mechanism}
The fixed link is $l_1=4\,\mathrm{cm}$, and the input link is
$l_2=9\,\mathrm{cm}$. Choose integer coupler and output lengths.
\begin{equation*}
  l_3=\blank\,\mathrm{cm},\qquad l_4=\blank\,\mathrm{cm}.
\end{equation*}
\prompt{Identify $S,L,P,Q$; demonstrate the strict inequality
$S+L<P+Q$; confirm that the fixed link $l_1$ is the shortest.}

\section{Part C: Numerical Position Analysis of a Planar Four-Bar}
Take $A^*=(0,0)$ and $B^*=(r_1,0)$. The input link runs from $A^*$
to $A$, the coupler from $A$ to $B$, and the output link from $B^*$ to $B$.
Angles $\theta_2,\theta_3,\theta_4$ are measured counterclockwise from
the positive $x$-axis. The loop closure is
\begin{equation}
  r_2 e^{i\theta_2}+r_3 e^{i\theta_3}
  =r_1+r_4 e^{i\theta_4}.
  \label{eq:loop}
\end{equation}
For numerical solution, use the two residual equations
\begin{align}
  f_1(\theta_3,\theta_4)
    &=r_2\cos\theta_2+r_3\cos\theta_3-r_1-r_4\cos\theta_4=0,\\
  f_2(\theta_3,\theta_4)
    &=r_2\sin\theta_2+r_3\sin\theta_3-r_4\sin\theta_4=0.
  \label{eq:residuals}
\end{align}
One suitable numerical method is Newton--Raphson iteration,
\begin{equation}
  \begin{bmatrix}\theta_3\\\theta_4\end{bmatrix}_{k+1}
  =\begin{bmatrix}\theta_3\\\theta_4\end{bmatrix}_{k}
   -J(\boldsymbol\theta_k)^{-1}
    \begin{bmatrix}f_1\\f_2\end{bmatrix}_{k},
  \qquad
  J=\begin{bmatrix}
    -r_3\sin\theta_3 & r_4\sin\theta_4\\
     r_3\cos\theta_3 &-r_4\cos\theta_4
  \end{bmatrix}.
  \label{eq:newton}
\end{equation}
Use radians during evaluation and convert the final angles to degrees.
State the initial guess and the selected assembly branch (open or crossed),
because two configurations may satisfy the same input data. After finding
the angles, compute the required diagonal $L=|A^*B|$ from
\begin{equation}
  L=\sqrt{(r_1+r_4\cos\theta_4)^2+(r_4\sin\theta_4)^2}.
  \label{eq:diagonal}
\end{equation}
Report a small final residual as a check of the numerical solution.

\subsection{Task 1}
Given $r_1=90$, $r_2=40$, $r_3=140$, $r_4=110$, and
$\theta_2=55^\circ$ (length units as supplied in the assignment).
\begin{align*}
  \text{Initial guess }(\theta_3,\theta_4)&=\prompt{angles and branch},\\
  \text{Numerical method and stopping criterion}&=\prompt{method and tolerance},\\
  \theta_3&=\blank^\circ,\qquad \theta_4=\blank^\circ,\\
  L=|A^*B|&=\blank,\qquad \lVert(f_1,f_2)\rVert=\blank.
\end{align*}
\prompt{Show the substituted residual equations and enough iteration or
solver detail to make the numerical result reproducible.}

\subsection{Task 2}
Given $r_1=70$, $r_2=55$, $r_3=120$, $r_4=95$, and
$\theta_2=115^\circ$ (length units as supplied in the assignment).
\begin{align*}
  \text{Initial guess }(\theta_3,\theta_4)&=\prompt{angles and branch},\\
  \text{Numerical method and stopping criterion}&=\prompt{method and tolerance},\\
  \theta_3&=\blank^\circ,\qquad \theta_4=\blank^\circ,\\
  L=|A^*B|&=\blank,\qquad \lVert(f_1,f_2)\rVert=\blank.
\end{align*}
\prompt{Show the substituted residual equations and enough iteration or
solver detail to make the numerical result reproducible.}

\clearpage
\section{Summary of Results}
The compact table below follows the grouping in the assignment's example.
Give the essential result in each cell; the detailed counts and calculations
remain in Parts A--C.

\begin{table}[H]
  \centering
  \caption{Answer summary for all tasks.}
  \small
  \begin{tabularx}{\linewidth}{@{}XXX@{}}
    \toprule
    Part A: mobility & Part B: classification and design & Part C: position\\
    \midrule
    Case 1: $M=\shortblank$ & Task 1, Case 1: \shortblank & Task 1: $L=\shortblank$, $\theta_3=\shortblank^\circ$, $\theta_4=\shortblank^\circ$\\
    Case 2: $M=\shortblank$ & Task 1, Case 2: \shortblank & Task 2: $L=\shortblank$, $\theta_3=\shortblank^\circ$, $\theta_4=\shortblank^\circ$\\
    Case 3: $M=\shortblank$ & Task 1, Case 3: \shortblank & \\
    & Task 2: $l_3=\shortblank\,\mathrm{cm}$, $l_4=\shortblank\,\mathrm{cm}$ & \\
    \bottomrule
  \end{tabularx}
\end{table}

% Final submission: replace student details and blue prompts, add three
% annotated Cinderella screenshots, and archive the PDF with three .cdy
% files as HW1_firstnameLASTNAME_xxxxxxxxx.zip.

\end{document}
